\newif\ifglobal

%%%%%%%%%%%%%%%%%%%%%%%
%%% Compile to view %%%
%%%%%%%%%%%%%%%%%%%%%%%

%%%%%%%%%%%%%%%%%%%%%%%%%%%%%%%%%%%%%%%%%%%%%%%%%%%
%%% Readme:                                     %%%
%%% https://www.overleaf.com/read/bwdjvfjksvjy/ %%%
%%%%%%%%%%%%%%%%%%%%%%%%%%%%%%%%%%%%%%%%%%%%%%%%%%%

% >>> M?: marks a module setting number or important notice

% >>> M4: Set {./relative-path} to external files for this module <<<
\def \modulefiles {./03-SYSMEM}
% To add an external file, use:
% (Replace '---' with actual filename)
% '\input{\modulefiles/tables/---.tex}' for tables
% '\includegraphics{\modulefiles/figures/---}' for figures
% '\subfile{\modulefiles/---}' for register files


% >>> M1: This module is in global project? <<<
% 'YES' - uncomment; 'NO' - comment out
\globaltrue

\ifglobal
    \documentclass[main\_\_CN.tex]{subfiles}

\else
    % Font "FandolSong-Regular" used by ctex does not have GSUB table.
    % It causes warnings that do not affect compilation and can be ignored.
    \PassOptionsToPackage{quiet}{fontspec}

    \documentclass[openany, 10pt]{book}

    % >>> M2: In ./00-shared/preamble-trm-module, also find
    %         settings for visibility of todo notes, watermark <<<
    \usepackage{preamble-trm-module}


    % Enable Chinese typesetting
    \usepackage{ctex}

    % >>> M3: Temporary labels in Chinese
    %         you can add your own labels there <<<
    \myexternaldocument{./00-shared/temp-labels-cn}

\fi %\ifglobal


\setcounter{secnumdepth}{4}


% >>> M5: If this module needs any updates in
% './00-shared/preamble-trm-module.sty'
% (include additional package, etc.), contact your module PM
% or create the file './00-shared/changelog.tex' make the
% necessary updates yourself and document them in this file <<<

% Make text left-aligned and allow latex to break pages naturally, comment out for full justification
\raggedright
\raggedbottom

\begin{document}


% Sets language into which pre-defined bits of text are translated
\selectlanguage{chinese}
\setchineseformatting

% Do not delete this block
\ifglobal
% Add the following front matter if
% this module is outside of global project
\else
    \InputIfFileExists{readme}{}{}
    \listoftodos
    {\let\clearpage\relax \listofdonetodos}
    \tableofcontents
    %\listoftables
    %\listoffigures
\fi



%%%%%%%%%%%%%%%%%%%%%%%%%
%%% TRM Module Begins %%%
%%%%%%%%%%%%%%%%%%%%%%%%%

\hypertarget{sysmem}{}
\chapter{系统和存储器} 
\label{mod:sysmem}


\section{概述}

ESP32-S2 采用哈佛结构 Xtensa$^®$ LX7 CPU 构成单核系统。所有的内部存储器、外部存储器以及外设都分布在 CPU 的总线上。

\section{主要特性}

\begin{itemize}
    \item {\bfseries 地址空间}
        \begin{itemize}
            \item 数据总线与指令总线总共有 4 GB（32 位）地址空间
    	    \item 464 KB 内部存储器指令地址空间
    	    \item 400 KB 内部存储器数据地址空间
    	    \item 1.77 MB  外设地址空间
    	    \item 7.5 MB 外部存储器指令虚地址空间
    	    \item 14.5 MB 外部存储器数据虚地址空间
    	    \item 320 KB 内部 DMA 地址空间
    	    \item 10.5 MB 外部 DMA 地址空间
        \end{itemize}
    \item {\bfseries 内部存储器}
        \begin{itemize}
            \item 128 KB Internal ROM
	        \item 320 KB Internal SRAM
	        \item 8 KB RTC FAST Memory
	        \item 8 KB RTC SLOW Memory
	    \end{itemize}
    \item {\bfseries 外部存储器}
        \begin{itemize}
            \item 最大支持 1 GB 外部 SPI flash
	        \item 最大支持 1 GB 外部 SPI RAM
	    \end{itemize}
    \item {\bfseries DMA}
        \begin{itemize}
            \item 9 个具有 DMA 功能的模块 / 外设
        \end{itemize}
\end{itemize}

图 \ref{figure:System_AddressMapping_Structure} 描述了系统结构与地址映射结构。

\begin{figure}[H]
    \centering
    \includegraphics[width=1.0\textwidth]{\modulefiles/figures/newchip_memorymapping_20191118}
    \caption{系统结构与地址映射结构}
    \label{figure:System_AddressMapping_Structure}
\end{figure}

\begin{tiplisting}
\vspace{-2em}
\begin{itemize}
    \item 图中灰色背景标注的地址空间不可用。
    \item 地址空间中可用的地址范围可能大于或小于实际可用的内存。
\end{itemize}

\end{tiplisting}

\section{功能描述}

\subsection{地址映射}

系统由一个哈佛结构 Xtensa$^®$ LX7 CPU 构成，CPU 具有 4 GB（32 位）的地址空间寻址能力。

地址 \hexprefix{}4000\_0000 以下的部分属于数据总线的地址范围，
地址 \hexprefix{}4000\_0000 \verb+~+ \hexprefix{}4FFF\_FFFF 部分为指令总线的地址范围，
地址 \hexprefix{}5000\_0000 及以上的部分是数据总线与指令总线共用的地址范围。

CPU 的数据总线与指令总线都为小端序。
CPU 可以通过数据总线进行字节、半字、字对齐的数据访问。
CPU 也可以通过指令总线进行数据访问，但必须是字对齐方式；非对齐数据访问会导致 CPU 工作异常。

CPU 能够：
\begin{itemize}
\item 通过数据总线与指令总线直接访问内部存储器；
\item 通过 cache 直接访问映射到地址空间的外部存储器；
\item 通过数据总线直接访问模块 / 外设。
\end{itemize}
表 \ref{table:Address_Mapping} 描述了 CPU 的数据总线与指令总线中的各段地址所能访问的目标。

系统中部分内部存储器与部分外部存储器既可以被数据总线访问也可以被指令总线访问，
这种情况下，CPU 可以通过多个地址访问到同一目标。

\subfile{\modulefiles/Address_Mapping_cn}

\subsection{内部存储器}\label{para:embedded memory}

内部存储器分为四个部分：Internal ROM (128 KB)、Internal SRAM (320 KB)、RTC FAST Memory (8 KB)、RTC SLOW Memory (8 KB)。

128 KB Internal ROM 分为 64 KB Internal ROM 0、64 KB Internal ROM 1 两部分。

320 KB Internal SRAM 分为 32 KB Internal SRAM 0、288 KB Internal SRAM 1 两部分。

RTC FAST Memory 与 RTC SLOW Memory 都为 SRAM。

表 \ref{table:Embedded_Memory_Address_Mapping} 列出了所有内部存储器以及访问内部存储器的
数据总线与指令总线地址段。

\subfile{\modulefiles/Embedded_Memory_Address_Mapping_cn.tex}



\subsubsection{Internal ROM 0}

Internal ROM 0 的容量为 64 KB，只读。如表 \ref{table:Embedded_Memory_Address_Mapping} 所示，CPU 只可以通过指令总线访问这部分存储器。


\subsubsection{Internal ROM 1}

Internal ROM 1的容量为64 KB，只读。如表 \ref{table:Embedded_Memory_Address_Mapping} 所示，CPU 可以通过数据或指令总线进行访问。

这两段地址同序访问 Internal ROM 1，即地址 \hexprefix{}3FFA\_0000 与 \hexprefix{}4001\_0000 访问到相同的字，\hexprefix{}3FFA\_0004 与 \hexprefix{}4001\_0004 访问到相同的字，\hexprefix{}3FFA\_0008 与 \hexprefix{}4001\_0008 访问到相同的字，以此类推。


\subsubsection{Internal SRAM 0}\label{Internal SRAM 0}

Internal SRAM 0 的容量为 32 KB，可读可写。如表 \ref{table:Embedded_Memory_Address_Mapping} 所示，CPU 可以通过数据或指令总线同序访问。


\hypertarget{Cache}{}
通过配置，这部分存储器中的 8 KB、16 KB、24 KB 或全部 32 KB 可以被 cache  占用，用来缓存外部存储器的数据。被 cache 占用的部分不可以被 CPU 访问，未被 cache 占用的部分仍然可以被 CPU 访问。




\subsubsection{Internal SRAM 1}

Internal SRAM 1 容量为 288 KB，可读可写。如表 \ref{table:Embedded_Memory_Address_Mapping} 所示，CPU 可以通过数据或指令总线同序访问。

\hypertarget{Trace}{}
Internal SRAM 1 存储器的 288 KB 地址空间由 18 个容量为 16 KB 的小存储器（子存储器）组成。其中至多有 1 个可以作为 Trace Memory 用于 CPU 的 Trace 功能，但 Trace Memory 无法被 CPU 访问。

\subsubsection{RTC FAST Memory}
\hypertarget{PRO_CPU_Only}{}
RTC FAST Memory 为 8 KB SRAM，可读可写。如表 \ref{table:Embedded_Memory_Address_Mapping} 所示，CPU 可以通过数据或指令总线同序访问。

\subsubsection{RTC SLOW Memory}

RTC SLOW Memory 为 8 KB SRAM，可读可写。如表 \ref{table:Embedded_Memory_Address_Mapping} 所示，CPU 可以通过数据／指令总线的共用地址段访问这部分存储器。

RTC SLOW Memory 也可以被当作一个外设来使用，此时 CPU 可以通过数据总线的地址段 \hexprefix{}3F42\_1000 \verb+~+ \hexprefix{}3F42\_2FFF 或地址段 \hexprefix{}6002\_1000 \verb+~+ \hexprefix{}6002\_2FFF 来访问它。








\subsection{外部存储器 } \label{para:external_Memory}

ESP32-S2 支持多个外部 QSPI / OSPI flash 和随机存储器 (RAM)。
ESP32-S2 还支持基于 XTS-AES 算法 的硬件加解密功能，从而保护开发者 flash 和 片外 RAM 中的程序和数据。

\subsubsection{外部存储器地址映射}

CPU 通过高速缓存 (cache) 来访问外部存储器。Cache 将根据 MMU 中的信息把 CPU 指令总线或数据总线 的地址变换为访问外部 flash 与片外 RAM 的实地址。经过变换的实地址所组成的实地址空间最大支持 1 GB 的外部 flash 与 1 GB 的片外 RAM。


通过高速缓存，ESP32-S2 一次最多可以同时有：
\begin{itemize}
    \item 7.5 MB 的指令总线 地址空间，通过指令缓存 (ICache) 以 64 KB 为单位 映射到 flash 或 片外 RAM，支持字节、半字、字对齐的读访问。 
    \item 4 MB 的只读数据总线 地址空间，通过指令缓存 (ICache) 以 64 KB 为单位 映射到 flash 或 片外 RAM， 支持 字节、半字、字对齐的读访问。
    \item 10.5 MB 的数据总线 地址空间，通过数据缓存 (DCache) 以 64 KB 为单位 映射到片外 RAM， 支持 字节、半字、字对齐的读写访问。 这部分地址空间 也可以用作只读数据空间，映射到 flash 与片外 RAM。
\end{itemize}


表 \ref{table:External_Memory_Address_Mapping} 列出了在访问 外部存储器时 CPU 的数据总线与指令总线 与 cache 的对应关系。

\subfile{\modulefiles/External_Memory_Address_Mapping_cn.tex}


\subsubsection{高速缓存}

如图 \ref{figure:Cache_arch} 所示， ESP32-S2 cache 采用分立结构，以便当 CPU 的指令总线和数据总线同时发起请求时，也可以迅速响应。 Cache 的存储空间与内部存储空间复用（详见章节 \ref{Internal SRAM 0}）。当 cache 缺失时，cache 控制器会向外部存储器发起请求，当 ICache 和 DCache 同时发起外部存储器请求时，由仲裁器决定谁先获得外部存储器的使用权。
ICache 和 DCache 的缓存大小可配置为 8 KB 和 16 KB， 块大小可配置为 16 B 和 32 B。


\begin{figure}[H]
    \centering
    \fbox{\includegraphics[width=0.45\textwidth]{\modulefiles/figures/Cache_arch_cn}}
     \caption{Cache 系统框图}
    \label{figure:Cache_arch}
\end{figure}


\subsubsection{Cache 操作}
ESP32-S2 cache 共有如下几种操作：
\begin{enumerate}
\item \textbf{失效} (Invalidate)：清除 Tag 的有效标志位 (Valid bit)。如果 CPU 接着去访问该数据，那么需要访问外部存储器。失效分为自动失效和手动失效。手动失效仅对 cache 中落入指定区域的地址对应的数据做失效处理，而自动失效会对 cache 中的所有数据做失效处理。ICache 和 DCache 均具有此功能。

\item \textbf{清除} (Clean)：清除 Tag 的脏标志位 (Dirty bit)，保留有效标志位。如果 CPU 接着去访问该数据，那么直接可以从 cache 中访问到。只有 DCache 具有此功能。

\item \textbf{写回} (Write-back)：清除 Tag 的脏块标志位，保留有效标志位， 并强制将对应地址的数据写回到外部存储器。如果 CPU 接着去访问该数据，那么可以直接从 cache 中访问到。只有 DCache 具有此功能。

\item \textbf{预取} (Preload)：Preload 功能用于将指令和数据提前加载到 cache 中。预取操作的最小单位为 1 个块。预取分为手动预取和自动预取，手动预取是指硬件按软件指定的虚地址预取一段连续的数据；自动预取是指硬件根据当前命中 / 缺失（取决于配置）的地址，自动预取一段连续的数据。

\item \textbf{锁定 / 解锁} (Lock / Unlock)：锁定分为预锁定和手动锁定，预锁定开启时，cache 在填充缺失数据到 cache memory 时将落在指定区域的数据锁定，手动锁定开启时，cache 检查已填充到 cache memory 的数据，并将落在指定区域的数据锁定。锁定区域的数据会一直保存在 cache 中，而不会被替换掉。但当所有路都被锁定时, cache 将进行正常替换，就像 cache 没有被锁定一样。解锁是锁定的逆操作。Cache 的手动失效操作、清除和写回操作均需要在解锁后才可使用。

\end{enumerate}



\subsection{ DMA 地址空间 }

ESP32-S2 可以借助直接存储访问 (direct memory access, DMA) 完成： \begin{itemize}
\item 模块 / 外设与内部存储器之间的数据搬运； \item 内部存储器与内部存储器之间的数据搬运；
\item 模块 / 外设与外部存储器之间的数据搬运；
\item 内部存储器与外部存储器之间的数据搬运。
\end{itemize}

DMA 可以通过与 CPU 数据总线完全相同的地址访问 Internal SRAM 0 与 Internal SRAM 1，
即 使用地址 \hexprefix{}3FFB\_0000 \verb+~+ \hexprefix{}3FFB\_7FFF 访问 Internal SRAM 0，
使用地址 \hexprefix{}3FFB\_8000 \verb+~+ \hexprefix{}3FFF\_FFFF 访问 Internal SRAM 1。
但  DMA 无法访问被 cache 占用的内部存储器。

另外， DMA 可以使用与 CPU 访问 DCache 相同的地址 (\hexprefix{}3F50\_0000 \verb+~+ \hexprefix{}3FF7\_FFFF) 来访问外部存储器，但只可以访问片外 RAM。 当 DCache 与 DMA 同时访问外部存储器时，务必要注意数据一致性问题。

系统中具有 DMA 功能的模块 / 外设总共有 9 个，如表 \ref{table:Module_with_DMA} 所示。其中，有些模块 / 外设通过 DMA 只可以访问内部存储器，有些既可以访问内部存储器又可以访问外部存储器。更多关于 DMA 的信息，请参考章节 \ref{mod:dma}：\textit{DMA 控制器}。

\subfile{\modulefiles/Module_with_DMA_cn.tex}



\subsection{模块 / 外设地址空间}

CPU 可以通过数据总线 \hexprefix{}3F40\_0000 \verb+~+ \hexprefix{}3F4F\_FFFF、数据 / 指令总线的共用地址段 \hexprefix{}6000\_0000 \verb+~+ \hexprefix{}600B\_FFFF 和 \hexprefix{}6180\_0000 \verb+~+ \hexprefix{}6180\_3FFF 来访问模块 / 外设。

\subsubsection{外设总线名称约定} \label{Convention: Peripheral Bus Name}
为了后续描述的方便，作如下约定：
\begin{itemize}
    \item \textbf{PeriBus1}\label{PeripheralBus1}：总线地址段 \hexprefix{}3F40\_0000 \verb+~+ \hexprefix{}3F4F\_FFFF 被记作“PeriBus1 总线”，简记为“\hyperref[PeripheralBus1]{PeriBus1}”，它的基地址为 \hexprefix{}3F40\_0000。
    \item \textbf{PeriBus2}\label{PeripheralBus2}：总线地址段 \hexprefix{}6000\_0000 \verb+~+ \hexprefix{}600B\_FFFF 和 \hexprefix{}6180\_0000 \verb+~+ \hexprefix{}6180\_3FFF 整体被记作“PeriBus2 总线”，简记为“\hyperref[PeripheralBus2]{PeriBus2}”，它的基地址为 \hexprefix{}6000\_0000。
\end{itemize}
在后续章节中出现的所有“\hyperref[PeripheralBus1]{PeriBus1}” 和“\hyperref[PeripheralBus2]{PeriBus2}” 都指代对应的总线地址段。


\subsubsection{外设总线区别}\label{differences_between_buses}

相比于 CPU 通过   \hyperref[PeripheralBus2]{PeriBus2}  访问模块 / 外设，CPU 通过  \hyperref[PeripheralBus1]{PeriBus1}  访问模块 / 外设效率更高。但是 \hyperref[PeripheralBus1]{PeriBus1} 有预测性读 (speculative read) 的特点，不能保证每一次的读访问都是真实有效的，因此，在访问模块 / 外设中的某些特殊寄存器如 FIFO 寄存器时，必须使用   \hyperref[PeripheralBus2]{PeriBus2}  进行访问。

另外，\hyperref[PeripheralBus1]{PeriBus1} 会打乱总线上的读写操作的先后顺序以提升性能， 这可能会导致对读写操作的先后顺序有严格要求的程序发生崩溃，对于这种情况，请在程序语句之前增加 volatile，也可以改用  \hyperref[PeripheralBus2]{PeriBus2}  进行访问。

\subsubsection{模块 / 外设地址空间列表}

表 \ref{table:Peripheral_Address_Mapping} 详细列出了模块 / 外设地址空间的各段地址与其能访问到的模块 / 外设的映射关系。请注意，“边界地址” 一栏中的数值是相对于总线基地址的地址偏移量，而非绝对地址。访问某一模块 / 外设的绝对地址等于总线基地址加上相应的地址偏移量。

\subfile{\modulefiles/Peripheral_Address_Mapping_cn.tex}

\subsubsection{\hyperref[PeripheralBus1]{PeriBus1} 访问受限地址列表} \label{special register inside peripherals}

如前文 \ref{differences_between_buses} 所述，\hyperref[PeripheralBus1]{PeriBus1} 有预测性读 (speculative read) 的特点，因此被禁止读访问模块 / 外设的 FIFO 寄存器。表 \ref{table:Special_Registers_for_Dport_Bus} 列出了禁止被 \hyperref[PeripheralBus1]{PeriBus1} 读访问的地址。另外，4 个用户自定义寄存器可以根据用户需求进行配置，用于将更多地址添加到访问受限地址列表中。更多信息请参考章节 \ref{mod:permctrl}：\textit{权限管理}。

\subfile{\modulefiles/Special_Registers_for_Dport_Bus_cn.tex}













\end{document}
